\chapter{ארכיטקטורת \textenglish{Transformer}}

\section{סקירה כללית}

ארכיטקטורת \textenglish{Transformer} מורכבת מחלקים עיקריים:
מקודד (\textenglish{encoder}), מפענח (\textenglish{decoder}), וחלק ליניארי לחיזוי תוצאה.
כל רכיב תורם לתפקיד משלו בקשת העיבוד.

\section{המקודד}

המקודד משנה את הקלט למטריצה של ייצוגים מזוקנים.
כל שכבת מקודד מורכבת משתי תת־שכבות:

\begin{enumerate}
\item תת־שכבת תשומת הלב מרובת ראשים (\textenglish{Multi-Head Attention})
\item רשת עצביות בעלת הנחיה קדימה (\textenglish{Feed-Forward Network})
\end{enumerate}

שתיהן מוקפות בתוך \textenglish{LayerNorm} ובחיבור בדלג (\textenglish{residual connection}).

\section{המפענח}

המפענח פועל בדומה למקודד, אך עם הבדל חשוב:
הוא מפעיל תשומת הלב "מוסכת" כדי שכל אסימון יכול רק להתייחס לאסימונים קודמים.

\section{תשומת הלב מרובת ראשים}

חישוב תשומת הלב מרובת ראשים מתבצע בנוסחה הבאה:

\begin{equation}
\text{MultiHead}(Q, K, V) = \text{Concat}(\text{head}_1, \ldots, \text{head}_h) W^O
\end{equation}

כאשר כל ראש מחושב כך:

\begin{LTR}
\begin{verbatim}
head_i = Attention(Q W_i^Q, K W_i^K, V W_i^V)
\end{verbatim}
\end{LTR}

\section{הרשת העצביות בעלת ההנחיה הקדימה}

כל משכבה יש רשת עצביות זעירה בעלת שתי שכבות ליניאריות:

\begin{equation}
\text{FFN}(x) = \max(0, xW_1 + b_1)W_2 + b_2
\end{equation}

זה נקרא גם בשם "תוספת בעלת שתי שכבות" ואו \textenglish{position-wise feed-forward network}.

\section{קידוד מיקום}

מכיוון ש־\textenglish{Transformer} עובד על כל הקלט בו־זמנית (בניגוד ל־\textenglish{RNN}),
המודל צריך כיצד לדעת את הסדר של האסימונים.
זה נעשה באמצעות קידוד מיקום המתווסף לייצוג הקלט:

\begin{equation}
PE_{(pos, 2i)} = \sin(pos / 10000^{2i/d_{model}})
\end{equation}

\begin{equation}
PE_{(pos, 2i+1)} = \cos(pos / 10000^{2i/d_{model}})
\end{equation}

\section{יישום \textenglish{PyTorch}}

דוגמה של יישום בסיסי בעזרת \textenglish{PyTorch}:

\begin{LTR}
\begin{verbatim}
import torch.nn as nn

class TransformerBlock(nn.Module):
    def __init__(self, d_model, num_heads, d_ff):
        super().__init__()
        self.attn = MultiHeadAttention(d_model, num_heads)
        self.ffn = FeedForward(d_model, d_ff)
        self.ln1 = nn.LayerNorm(d_model)
        self.ln2 = nn.LayerNorm(d_model)
    
    def forward(self, x):
        x = x + self.attn(self.ln1(x))
        x = x + self.ffn(self.ln2(x))
        return x
\end{verbatim}
\end{LTR}

קוד זה מראה את המבנה הבסיסי של בלוק \textenglish{Transformer} יחיד.
